\section{Model, reduction, and the quantities we design}\label{sec:model}

\subsection{Laminate constitutive model}\label{sec:clt}
A unidirectional CFRP lamina ($E_1=\SI{131}{GPa}$, $E_2=\SI{9}{GPa}$,
$G_{12}=\SI{5}{GPa}$, $\nu_{12}=0.27$, $\rho=\SI{1600}{kg/m^3}$) has the on-axis
plane-stress stiffness $\bm Q$, with $Q_{11}=E_1/(1-\nu_{12}\nu_{21})$,
$Q_{22}=E_2/(1-\nu_{12}\nu_{21})$, $Q_{12}=\nu_{12}Q_{22}$, $Q_{66}=G_{12}$ and
$\nu_{21}=\nu_{12}E_2/E_1$. In engineering Voigt notation (strain
$\bm s=[s_{xx},s_{yy},2s_{xy}]^{\!\top}$), rotating the fibers through $\theta$ gives
\begin{equation}
\bar{\bm Q}(\theta)=\bm T(\theta)\,\bm Q\,\bm T(\theta)^{\!\top},\qquad
\bm T(\theta)=\begin{bmatrix}c^2&s^2&-2sc\\ s^2&c^2&2sc\\ sc&-sc&c^2-s^2\end{bmatrix},
\quad (c,s)=(\cos\theta,\sin\theta).
\label{eq:rot}
\end{equation}
Its entries are trigonometric polynomials of degree four in $\theta$ (Tsai--Pagano \cite{tsai1968}):
\begin{equation}
\bar{\bm Q}(\theta)=\bm\Gamma_0+\cos2\theta\,\bm\Gamma_1+\sin2\theta\,\bm\Gamma_2
+\cos4\theta\,\bm\Gamma_3+\sin4\theta\,\bm\Gamma_4,
\label{eq:qbar}
\end{equation}
\begin{gather*}
\bm\Gamma_0=\begin{bmatrix}U_1&U_4&0\\U_4&U_1&0\\0&0&U_5\end{bmatrix},\quad
\bm\Gamma_1=\begin{bmatrix}U_2&0&0\\0&-U_2&0\\0&0&0\end{bmatrix},\quad
\bm\Gamma_2=\tfrac{U_2}{2}\begin{bmatrix}0&0&1\\0&0&1\\1&1&0\end{bmatrix},\\
\bm\Gamma_3=U_3\begin{bmatrix}1&-1&0\\-1&1&0\\0&0&-1\end{bmatrix},\quad
\bm\Gamma_4=U_3\begin{bmatrix}0&0&1\\0&0&-1\\1&-1&0\end{bmatrix},
\end{gather*}
with the material invariants
$U_1=\tfrac18(3Q_{11}+3Q_{22}+2Q_{12}+4Q_{66})$, $U_2=\tfrac12(Q_{11}-Q_{22})$,
$U_3=\tfrac18(Q_{11}+Q_{22}-2Q_{12}-4Q_{66})$, $U_4=\tfrac18(Q_{11}+Q_{22}+6Q_{12}-4Q_{66})$,
$U_5=\tfrac18(Q_{11}+Q_{22}-2Q_{12}+4Q_{66})$. $\bm\Gamma_0$ is isotropic. All the
dependence on orientation is carried by $U_2$ and $U_3$, which vanish for an isotropic
material. For the CFRP above, $U_2/U_1=1.10$ and $U_3/U_1=0.26$.

A symmetric laminate ($\bm B=\bm 0$) of thickness $h$ has the membrane and bending
stiffnesses
\begin{equation}
\bm A=h\Big[\bm\Gamma_0+\textstyle\sum_{i=1}^4\xi^A_i\bm\Gamma_i\Big],\qquad
\bm D=\tfrac{h^3}{12}\Big[\bm\Gamma_0+\textstyle\sum_{i=1}^4\xi^D_i\bm\Gamma_i\Big],
\label{eq:lampar}
\end{equation}
with the lamination parameters \cite{tsai1968,abdalla2007}
\begin{equation}
\bm\xi^A=\int_{-1/2}^{1/2}\bm v(\theta(\bar z))\,\dd\bar z,\qquad
\bm\xi^D=12\int_{-1/2}^{1/2}\bar z^2\,\bm v(\theta(\bar z))\,\dd\bar z,\qquad
\bm v(\theta)=(\cos2\theta,\sin2\theta,\cos4\theta,\sin4\theta),
\label{eq:xi}
\end{equation}
where $\theta(\bar z)$ is the ply angle at normalized height $\bar z=z/h$. Two facts
organize everything that follows.
\begin{itemize}[itemsep=1pt]
\item A \emph{single-ply-like} laminate (one angle through the thickness) has
$\bm\xi^A=\bm\xi^D=\bm v(\theta)$, so $\bm D=\tfrac{h^2}{12}\bm A$ pointwise.
\item A \emph{stacking sequence} makes $\bm\xi^D\ne\bm\xi^A$. Because $\bm\xi^D$ weights
the outer plies by $12\bar z^2$, the stacking sequence separates bending from membrane
stiffness \emph{through the thickness}.
\end{itemize}
A \emph{steered} laminate makes $\bm\xi^{A,D}$ functions of the in-plane position
$\bm x$.

\subsection{The laminated von K\'arm\'an plate}\label{sec:plate}
A rectangular plate $\Omega=[0,a]\times[0,b]$ ($a=\SI{0.20}{m}$, $b=\SI{0.15}{m}$,
$h=\SI{2}{mm}$) carries the transverse deflection $w$ and in-plane displacements
$\bm u=(u,v)$. The von K\'arm\'an membrane strain and curvature are
\begin{equation}
\bm\varepsilon=\bm\varepsilon_L(\bm u)+\bm\varepsilon_{NL}(w),\quad
\bm\varepsilon_L=\begin{bmatrix}u_{,x}\\v_{,y}\\u_{,y}+v_{,x}\end{bmatrix},\quad
\bm\varepsilon_{NL}=\tfrac12\begin{bmatrix}w_{,x}^2\\w_{,y}^2\\2w_{,x}w_{,y}\end{bmatrix},\qquad
\bm\kappa=-\begin{bmatrix}w_{,xx}\\w_{,yy}\\2w_{,xy}\end{bmatrix},
\label{eq:kin}
\end{equation}
and, with $\bm A(\bm x),\bm D(\bm x)$ from \eqref{eq:lampar}, the strain energy is
\begin{equation}
U=\tfrac12\!\int_\Omega\!\bm\kappa^{\!\top}\bm D\,\bm\kappa\,\dd\Omega
+\tfrac12\!\int_\Omega\!(\bm\varepsilon_L+\bm\varepsilon_{NL})^{\!\top}\bm A\,
(\bm\varepsilon_L+\bm\varepsilon_{NL})\,\dd\Omega .
\label{eq:energy}
\end{equation}
With the kinetic energy $T=\tfrac12\int\rho h(\dot w^2+|\dot{\bm u}|^2)$, Hamilton's
principle gives the laminated F\"oppl--von K\'arm\'an equations \cite{amabili2008},
\begin{equation}
N_{\alpha\beta,\beta}=\rho h\,\ddot u_\alpha,\qquad
M_{\alpha\beta,\alpha\beta}+\big(N_{\alpha\beta}w_{,\beta}\big)_{,\alpha}+p=\rho h\,\ddot w,
\qquad \bm N=\bm A(\bm\varepsilon_L+\bm\varepsilon_{NL}),\ \bm M=\bm D\bm\kappa .
\label{eq:fvk}
\end{equation}
The edges are simply supported and immovable ($w=0$, $\bm u=\bm 0$ on $\partial\Omega$).
We discretize by Galerkin's method on sine bases, which satisfy these conditions term by
term:
\begin{equation}
w=\sum_{k=1}^{N_w}q_k\,\phi_k,\quad \phi_k=\sin\tfrac{m_k\pi x}{a}\sin\tfrac{n_k\pi y}{b};
\qquad
u,v=\sum_{j=1}^{N_s}s^{u,v}_j\,\psi_j,\quad \psi_j=\sin\tfrac{i_j\pi x}{a}\sin\tfrac{l_j\pi y}{b}.
\end{equation}
We use $N_w=64$ ($m,n\le8$) and $N_s=36$ per component. For the robustness study of
Sec.~\ref{sec:robust} the same evaluator is built on two other bases: clamped edges
($w=\partial_nw=0$) with $\phi_k=\sin\tfrac{\pi x}{a}\sin\tfrac{m_k\pi x}{a}\,
\sin\tfrac{\pi y}{b}\sin\tfrac{n_k\pi y}{b}$, which reproduces the clamped isotropic square
plate's frequency parameter $35.99$ to $0.15\%$; and in-plane traction-free (``movable'') edges,
with a Chebyshev polynomial basis of degree $\le10$ for $u,v$, converged to four digits in
$\gamma$. The bases are converged
for the steered designs: Table~\ref{tab:conv} shows that the steered relaxation needs at least
five to six in-plane modes per direction; the three-mode basis that suffices for a uniform
laminate under-resolves it badly ($J/J_{\UD}=\nConvThree$ instead of $\nConvBase$).

\begin{table}[!tbp]
\centering\footnotesize
\caption{Convergence of the evaluator on the reported held-frequency steered design
(generated by \texttt{conv\_table.py}). $J/J_{\UD}$ and the relaxed share are converged on
the bold basis; three in-plane modes per direction under-resolve the steered relaxation.
The steered frequency is a Ritz upper bound that falls by $\nConvFreqDrop\%$ as the transverse basis
grows (the UD frequency does not change), so ``held frequency'' holds to about $\nConvFreqDrop\%$
across bases and $\nHeldTol\%$ against the independent FE model.
Bold: the basis used throughout.}
\label{tab:conv}
\IfFileExists{conv_table.tex}{\begin{tabular}{@{}ccccccc@{}}
\toprule
transverse $m,n\le$ & in-plane $i,j\le$ & $f_{1,\UD}$ (Hz) & $f_1$ (Hz) & $J/J_{\UD}$ & relaxed share (UD) & relaxed share (steered)\\ \midrule
8 & 3 & 256.3 & 256.3 & 0.523 & 33\% & 31\%\\
8 & 4 & 256.3 & 256.3 & 0.420 & 33\% & 44\%\\
8 & 5 & 256.3 & 256.3 & 0.375 & 33\% & 50\%\\
\bf 8 & \bf 6 & \bf 256.3 & \bf 256.3 & \bf 0.366 & \bf 33\% & \bf 51\%\\
8 & 8 & 256.3 & 256.3 & 0.360 & 33\% & 52\%\\
12 & 6 & 256.3 & 254.6 & 0.370 & 33\% & 51\%\\
12 & 8 & 256.3 & 254.6 & 0.364 & 33\% & 52\%\\
16 & 8 & 256.3 & 253.9 & 0.366 & 33\% & 51\%\\
\bottomrule
\end{tabular}
}{\textit{table pending}}
\end{table}

With $\bm x=[\bm q;\bm s]$ the system is
\begin{equation}
\bm M\ddot{\bm x}+\bm C\dot{\bm x}+\bm K\bm x+\bm B_2(\bm x,\bm x)+\bm T_3(\bm x,\bm x,\bm x)=\bm f(t),
\label{eq:eom}
\end{equation}
with
$M_{kl}=\rho h\int\phi_k\phi_l$ (likewise for $\psi$),
$\bm K=\mathrm{blkdiag}(\bm K_b,\bm K_{\mathrm{in}})$,
$(\bm K_b)_{kl}=\int\bm\kappa_k^{\!\top}\bm D\bm\kappa_l$ and
$(\bm K_{\mathrm{in}})_{jl}=\int\bm\varepsilon_{L,j}^{\!\top}\bm A\bm\varepsilon_{L,l}$.
$\bm B_2$ is the gradient of $\int\bm\varepsilon_L^{\!\top}\bm A\bm\varepsilon_{NL}$ (the
$\bm u\leftrightarrow w^2$ coupling) and $\bm T_3$ is the gradient of
$\tfrac12\int\bm\varepsilon_{NL}^{\!\top}\bm A\bm\varepsilon_{NL}$ (the membrane cubic).
Because $\bm B=\bm 0$, the \emph{linear} problem decouples into a bending block and an
in-plane block. The mass is design-independent: every design in this paper has the
same plate, the same material and the same thickness.

The \emph{beam} is the one-dimensional specialization: a clamped--clamped strip
($L=\SI{0.30}{m}$, width $\SI{20}{mm}$, $h=\SI{2}{mm}$) with axial $u$ and transverse
$w$, $\varepsilon=u'+\tfrac12w'^2$ and $\kappa=-w''$. For a single-ply laminate,
$A=bhE(x)$ and $D=\tfrac{bh^3}{12}E(x)$ with $E(x)=\bar Q_{11}(\theta(x))$. It is
discretized by Hermite-cubic/linear elements, and the axial ends are immovable.

\subsection{Spectral submanifold and backbone}\label{sec:ssm}
In first-order form $\bm z=[\bm x;\dot{\bm x}]$,
$\bm{\mathcal B}\dot{\bm z}=\bm{\mathcal A}\bm z+\bm F(\bm z)$ with
$\bm{\mathcal B}=\begin{bsmallmatrix}\bm C&\bm M\\\bm M&\bm0\end{bsmallmatrix}$,
$\bm{\mathcal A}=\begin{bsmallmatrix}-\bm K&\bm0\\\bm0&\bm M\end{bsmallmatrix}$ and
$\bm F=[-\bm B_2-\bm T_3;\bm0]$. Let $\lambda_1,\bar\lambda_1$ be the eigenvalues of the
fundamental, with right eigenvector $\bm v_1$ and left eigenvector $\bm u_1$
($\bm u_1^*\bm{\mathcal A}=\lambda_1\bm u_1^*\bm{\mathcal B}$,
$\bm u_1^*\bm{\mathcal B}\bm v_1=1$). The SSM $\mathcal W$ is the unique smoothest
invariant manifold tangent to their eigenspace \cite{haller2016}. It exists if
$\lambda_1$ is the slowest pair and no outer eigenvalue equals
$m_1\lambda_1+m_2\bar\lambda_1$ up to the truncation order. Parametrizing
$\bm z=\bm W(\bm p)$ and $\dot{\bm p}=\bm R(\bm p)$ gives the invariance equation
$\bm{\mathcal B}\,\mathrm D\bm W\,\bm R=\bm{\mathcal A}\bm W+\bm F(\bm W)$. Expanding both
maps in monomials $\bm p^{\bm m}$, it separates order by order into the homological
equations
\begin{equation}
(\Lambda_{\bm m}\bm{\mathcal B}-\bm{\mathcal A})\bm w_{\bm m}
+\bm{\mathcal B}(\bm v_1R^1_{\bm m}+\bar{\bm v}_1R^2_{\bm m})=\bm G_{\bm m},
\qquad \Lambda_{\bm m}=m_1\lambda_1+m_2\bar\lambda_1,
\label{eq:homological}
\end{equation}
where $\bm G_{\bm m}$ collects lower-order terms (Appendix~\ref{app:coeffs}). At the
near-resonant indices $\bm m=(k{+}1,k)$ the operator is (near-)singular. Projecting onto
$\bm u_1$ gives the reduced (normal-form) coefficients
$\gamma_{2k+1}=\bm u_1^*\bm G_{(k+1,k)}$. On the orbit
$p_1=\tfrac\rho2e^{i\vartheta}$, $p_2=\bar p_1$, the backbone and the modal decay rate are
\begin{equation}
\Om(\rho)=\Imag\Big(\lambda_1+\sum_{k\ge1}\gamma_{2k+1}\big(\tfrac{\rho}{2}\big)^{2k}\Big),\qquad
\alpha(\rho)=\Real(\cdots).
\label{eq:backbone}
\end{equation}
The implementation is verified against the analytic Duffing backbone through fifth
order ($3/8$ and $-21/256$), against full-order EPMC continuation, and against two
published results \cite{pozzi2025} (Appendix~\ref{app:calib}). The master
eigenvector is computed by Rayleigh-quotient inverse iteration on the complex-symmetric
pencil, with $\bm u_1=\bar{\bm v}_1$. A dense QZ solve of the stiff first-order pencil was
found to lose up to ten digits (eigenvector condition number $\sim10^{14}$ for laminated
beams) and is no longer used for symmetric systems (regression test in
\texttt{tests/test\_relaxed\_plate.py}).

\subsection{The relaxed hardening coefficient}\label{sec:relaxed}
Let the fundamental be the lowest Ritz pair of the bending block,
$\bm K_b\bm v=\omega_1^2\bm M_w\bm v$ with $\bm v^{\!\top}\bm M_w\bm v=1$. Its shape
$w_1=\sum_kv_k\phi_k$ \emph{adapts} to the stiffness field. This is the exact linear
master mode of \eqref{eq:eom}, because the blocks decouple when $\bm B=\bm0$. Write its
curvature as $\bm\kappa_1$ and its stretch as
\begin{equation}
\bm m=\big[w_{1,x}^2,\;w_{1,y}^2,\;2w_{1,x}w_{1,y}\big]^{\!\top}=2\bm\varepsilon_{NL}(w_1).
\end{equation}

\begin{proposition}[Relaxed cubic, variational form]\label{prop:relaxed}
Let the in-plane field be slaved quasi-statically to $q^2$, the square of the
fundamental's modal amplitude. Then the cubic coefficient of the single-mode reduced
dynamics $\ddot q+\omega_1^2q+\gamma q^3=0$ is
\begin{equation}
\gamma=\min_{\bm s}\;\tfrac12\int_\Omega(\bm m-2\bm\varepsilon_L(\bm s))^{\!\top}\bm A\,
(\bm m-2\bm\varepsilon_L(\bm s))\,\dd\Omega
=k_{nl}-2\,\bm g^{\!\top}\bm K_{\mathrm{in}}^{-1}\bm g
=\tfrac12\int_\Omega\bm\varepsilon_r^{\!\top}\bm A\,\bm\varepsilon_r\,\dd\Omega,
\label{eq:gamma}
\end{equation}
with $k_{nl}=\tfrac12\int\bm m^{\!\top}\bm A\bm m$,
$\bm g=\tfrac12\int\bm\varepsilon_L(\cdot)^{\!\top}\bm A\bm m$,
$\bm s^\star=\bm K_{\mathrm{in}}^{-1}\bm g$ and the relaxed stretch
$\bm\varepsilon_r=\bm m-2\bm\varepsilon_L(\bm s^\star)$.
\end{proposition}
\begin{proof}
With $w=qw_1$ and $\bm u=-q^2\bm s$, the membrane strain is
$\bm\varepsilon=\tfrac{q^2}{2}(\bm m-2\bm\varepsilon_L(\bm s))$. The membrane energy is
therefore $\tfrac{q^4}{8}\int(\bm m-2\bm\varepsilon_L(\bm s))^{\!\top}\bm A(\cdot)
=\tfrac{q^4}{4}F(\bm s)$, where $F(\bm s)=k_{nl}-4\bm s^{\!\top}\bm g+2\bm s^{\!\top}
\bm K_{\mathrm{in}}\bm s$. Static equilibrium of the in-plane field minimizes $F$, at
$\bm s^\star=\bm K_{\mathrm{in}}^{-1}\bm g$, with $F(\bm s^\star)=k_{nl}-2\bm g^{\!\top}
\bm K_{\mathrm{in}}^{-1}\bm g$. The reduced potential $\tfrac{\gamma}{4}q^4$ then gives
\eqref{eq:gamma}. The last form is $F(\bm s^\star)$ written out.
\end{proof}

Proposition~\ref{prop:relaxed} has three uses. (i) $k_{nl}$ is the restrained
(immovable in-plane field) upper bound, and $2\bm g^{\!\top}\bm K_{\mathrm{in}}^{-1}\bm g$
is the \emph{relaxation}. Its share of $k_{nl}$ is $\nRelaxUD\%$ for the unidirectional
plate. (ii) Because $\gamma$ is a \emph{minimum} over $\bm s$, its design derivative
needs no in-plane adjoint (envelope theorem; Sec.~\ref{sec:sens}). (iii) It writes the
hardening as the integral of an energy density (Sec.~\ref{sec:theory}).

\begin{proposition}[Agreement with the SSM]\label{prop:ssmcond}
For an undamped system \eqref{eq:eom} with $\bm B=\bm 0$, the cubic SSM coefficient of the
fundamental satisfies
$\Imag\gamma_3=\tfrac{3}{2\omega_1}\gamma\,[1+O(4\omega_1^2/\omega_{\mathrm{in}}^2)]$,
where $\omega_{\mathrm{in}}$ is the lowest in-plane frequency.
\end{proposition}
\noindent The proof is in Appendix~\ref{app:ssmcond}. In short: the $(1,1)$ in-plane
homological solve is exactly static ($\Lambda_{(1,1)}=0$), and the $(2,0)$ solve is
static up to the resolvent factor
$(\bm K_{\mathrm{in}}-4\omega_1^2\bm M_{\mathrm{in}})^{-1}$. For our plate
$4\omega_1^2/\omega_{\mathrm{in}}^2=\nRhoIn$ ($f_{\mathrm{in}}\approx\nFin$\,kHz on the full multi-mode plate), and \eqref{eq:gamma} reproduces the full
multi-mode SSM coefficient to within $\nSSMvsEval\%$ (Sec.~\ref{sec:confirm}).

\subsection{The backbone in physical amplitude, and the design quantity $J$}
For $\ddot q+\omega_1^2q+\gamma q^3=0$ the cubic SSM gives
$\Om=\omega_1+\tfrac38(\gamma/\omega_1)q^2$. The mid-plate deflection of the mode is
$a=w_cq$, with $w_c=w_1(a/2,b/2)$, so
\begin{equation}
\frac{\Om}{\omega_1}=1+\frac38J\,a^2=1+\beta\Big(\frac{a}{h}\Big)^2,\qquad
J=\frac{\gamma}{\omega_1^2w_c^2},\qquad \beta=\tfrac38Jh^2 .
\label{eq:J}
\end{equation}
$J$ (units m$^{-2}$) is the \emph{frequency-normalized detuning}. It is invariant to the
mode normalization, and it is what we design. The unidirectional plate has
$f_1=\nUDf$\,Hz, $\beta=\nEvBetaUD$ and a $\nUDdet\%$ shift at $a/h=0.4$.
The cubic truncation reproduces the full-order EPMC detuning to within $\nEPMCAgree\%$ over the
amplitudes EPMC reached, $a/h\le\nEPMCmaxUD$ (UD) and $\le\nEPMCmaxOpt$ (steered) (Appendix~\ref{app:calib}).

\subsection{Forced response and the jump-free amplitude}\label{sec:jumpth}
With harmonic forcing $\bm f\cos\Om t$ and damping $\bm C=2\zeta\omega_1\bm M$, seeking
$p_1=\tfrac\rho2e^{i\Om t}$ on the SSM gives the leading-order forced-response curve
\begin{equation}
\big[(\Om-\Om(\rho))^2+\alpha(\rho)^2\big]\rho^2=|f_1|^2,\qquad f_1=\bm u_1^*[\bm f;\bm0].
\label{eq:frc}
\end{equation}

\begin{proposition}[Jump-free amplitude]\label{prop:cusp}
For the cubic backbone \eqref{eq:J} with $\alpha=\zeta\omega_1$, the forced response
\eqref{eq:frc} is single-valued in $\Om$ for every drive whose peak (resonant) amplitude is
below
\begin{equation}
\frac{a_c}{h}=\Big(\frac{8\zeta}{3\sqrt3\,\beta}\Big)^{1/2}\;\propto\;J^{-1/2},
\label{eq:ac}
\end{equation}
and multivalued (a fold pair with jump and hysteresis) above it. At the critical drive the
vertical tangent sits at the lower amplitude $(\sqrt3/2)\,a_c$ on the response curve.
\end{proposition}
\begin{proof}
Write $x=a^2$, $c=\omega_1\beta/h^2$, $\mu=\zeta\omega_1$, and let $F$ be the drive
expressed in physical amplitude. Then \eqref{eq:frc} reads
$[(\sigma-cx)^2+\mu^2]x=F^2$, with $\sigma=\Om-\omega_1$. Put $\sigma=cx+s$. The curve is
$x=F^2/(s^2+\mu^2)$, $\sigma=cF^2/(s^2+\mu^2)+s$, and it folds where $\dd\sigma/\dd s=0$,
i.e.\ where $2cF^2s=(s^2+\mu^2)^2$. A fold pair first appears when this equation has a
double root, i.e.\ when $2cF^2=4s(s^2+\mu^2)$ also holds. Combining the two conditions
gives $s^2+\mu^2=4s^2$, so $s=\mu/\sqrt3$, and the critical drive is
$F_c^2=(s^2+\mu^2)^2/(2cs)=8\mu^3/(3\sqrt3\,c)$. The vertical tangent lies at
$x=F_c^2/(s^2+\mu^2)=2\mu/(\sqrt3c)$, while the peak of the response (at $s=0$) is
$x_c=F_c^2/\mu^2=8\mu/(3\sqrt3\,c)$, which is \eqref{eq:ac}. Below $F_c$ the curve has no
fold; above it, it has two.
\end{proof}
\noindent Proposition~\ref{prop:cusp} is the consequence we test. Holding $\omega_1$,
$\zeta$ and $h$ and reducing $J$ by a factor $r$ widens the jump-free amplitude by
$r^{-1/2}$. In Sec.~\ref{sec:confirm} we compute $a_c$ \emph{numerically}, as the drive
at which the full-order damped SSM response first becomes multivalued, and compare it
with \eqref{eq:ac}.
